% !TEX program = pdflatex
% =====================================================================
%  CARNET D'ENTRAINEMENT — FICHE 7 (Loi des gaz)
%  THEME 1 — Lois historiques et transformations d'un gaz
%  TUTO
% =====================================================================
\documentclass[11pt,a4paper]{article}
\usepackage[utf8]{inputenc}
\usepackage[T1]{fontenc}
\usepackage{lmodern}
\usepackage[french]{babel}
\usepackage{amsmath,amssymb,mathtools}
\usepackage{icomma}
\usepackage{siunitx}
\sisetup{output-decimal-marker={,},per-mode=symbol}
\DeclareSIUnit{\atm}{atm}
\DeclareSIUnit{\bar}{bar}
\DeclareSIUnit{\mmHg}{mmHg}
\usepackage[version=4]{mhchem}
\usepackage{xcolor}
\usepackage{colortbl}
\usepackage{tikz}
\usepackage{pgfplots}
\usepackage[most]{tcolorbox}
\usepackage{enumitem}
\usepackage{array}
\usepackage{booktabs}
\usepackage{fancyhdr}
\usepackage[a4paper,margin=2.1cm,headheight=26pt,headsep=10pt]{geometry}
\usepackage{titlesec}
\usetikzlibrary{arrows.meta,positioning,calc,shapes.geometric,fit,backgrounds}
\pgfplotsset{compat=1.18}

\definecolor{primary}{RGB}{150,98,162}
\definecolor{primarydark}{RGB}{92,52,110}
\definecolor{defcol}{RGB}{0,114,178}
\definecolor{formulacol}{RGB}{198,93,9}
\definecolor{attncol}{RGB}{200,40,55}
\definecolor{excol}{RGB}{0,135,90}
\definecolor{methcol}{RGB}{90,90,100}
\definecolor{pcol}{RGB}{200,40,55}
\definecolor{vcol}{RGB}{0,135,90}
\definecolor{tcol}{RGB}{198,93,9}
\definecolor{ncol}{RGB}{120,94,200}
\definecolor{gazcol}{RGB}{0,120,190}

\tcbset{boxbase/.style={enhanced,breakable,boxrule=0.9pt,arc=2.5pt,
   left=3mm,right=3mm,top=2.2mm,bottom=2.2mm,fonttitle=\bfseries,coltitle=white}}
\newtcolorbox{idee}[1][]{boxbase,colback=defcol!6,colframe=defcol,title={#1}}
\newtcolorbox{formule}[1][]{boxbase,colback=formulacol!8,colframe=formulacol,title={#1}}
\newtcolorbox{piege}[1][]{boxbase,colback=attncol!7,colframe=attncol,title={#1}}
\newtcolorbox{methode}[1][]{boxbase,colback=methcol!8,colframe=methcol,sharp corners,title={#1}}

\pagestyle{fancy}\fancyhf{}
\fancyhead[L]{\small\textcolor{primarydark}{\textbf{Fiche 7} — Loi des gaz}}
\fancyhead[R]{\small\textcolor{primarydark}{\textsc{Thème 1 · Tuto}}}
\fancyfoot[C]{\small\thepage}
\renewcommand{\headrulewidth}{0.8pt}
\renewcommand{\headrule}{\hbox to\headwidth{\color{primary}\leaders\hrule height \headrulewidth\hfill}}
\titleformat{\section}{\large\bfseries\color{primarydark}}{\thesection}{0.6em}{}
\setlength{\parindent}{0pt}\setlength{\parskip}{4pt}

\begin{document}

\begin{center}
{\Large\bfseries\color{primarydark} Thème 1 — Lois historiques et transformations d'un gaz}\\[2pt]
{\color{primary}\itshape Tuto à lire avant de commencer les exercices}
\end{center}
\vspace{-2pt}
\textcolor{primary}{\hrule height 1pt}
\vspace{6pt}

\section*{1. Décrire un gaz : quatre nombres, et surtout leurs unités}

L'état d'un gaz est fixé par \textbf{quatre variables} : la pression \textcolor{pcol}{$p$}, le volume
\textcolor{vcol}{$V$}, la quantité de matière \textcolor{ncol}{$n$} et la température
\textcolor{tcol}{$T$}. Dans ce thème, \textbf{on ne fabrique ni ne détruit de gaz} : $n$ reste constant.
On étudie seulement comment $p$, $V$ et $T$ se répondent quand on comprime, chauffe ou refroidit.

\begin{piege}[La règle d'or : la température se met TOUJOURS en kelvins]
Dans toutes les lois des gaz, $T$ s'exprime en \textbf{kelvins} :
\[ \boxed{\,T(\si{\kelvin}) = \theta(\si{\degreeCelsius}) + 273\,} \]
Une transformation « à $\qty{27}{\degreeCelsius}$ » se calcule avec $T=\qty{300}{\kelvin}$, jamais avec $27$.
Diviser deux températures en \si{\degreeCelsius} (par ex.\ $87/27$) donne un résultat \emph{faux}.
\end{piege}

\textbf{Conversions de pression à connaître} (l'unité SI est le pascal, très petite) :
\begin{center}\small
$\qty{1}{\atm}=\qty{101325}{\pascal}\approx\qty{101.3}{\kilo\pascal}=\qty{760}{\mmHg}=\qty{1.013}{\bar}$
\qquad ; \qquad $\qty{1}{\bar}=\qty{e5}{\pascal}$ \qquad ; \qquad $\qty{1}{\meter\cubed}=\qty{1000}{\litre}$.
\end{center}

\section*{2. Les trois lois historiques : on fige deux variables, les deux autres se répondent}

Chaque loi correspond à une expérience où l'on \textbf{bloque deux variables}. La quantité $n$ étant
toujours constante ici, il suffit de repérer \emph{la troisième variable qui ne bouge pas}.

\begin{center}\renewcommand{\arraystretch}{1.35}\small
\begin{tabular}{l l l l}
\toprule
\rowcolor{primary!12}
\textbf{Nom} & \textbf{Ce qui est constant} & \textbf{Relation} & \textbf{En clair}\\
\midrule
Boyle-Mariotte      & $T$ (\emph{isotherme}) & $p_1V_1=p_2V_2$ & $p$ et $V$ en sens \emph{inverse}\\
Charles--Gay-Lussac & $p$ (\emph{isobare})   & $\dfrac{V_1}{T_1}=\dfrac{V_2}{T_2}$ & $V$ suit $T$ (kelvins)\\
Loi isochore        & $V$ (\emph{isochore})  & $\dfrac{p_1}{T_1}=\dfrac{p_2}{T_2}$ & $p$ suit $T$ (kelvins)\\
Avogadro            & $p$ et $T$             & $\dfrac{V_1}{n_1}=\dfrac{V_2}{n_2}$ & $V$ suit la quantité $n$\\
\bottomrule
\end{tabular}
\end{center}

\begin{formule}[La loi combinée : quand $p$, $V$ ET $T$ changent en même temps ($n$ fixe)]
\[ \boxed{\;\dfrac{p_1\,V_1}{T_1}=\dfrac{p_2\,V_2}{T_2}\;}\qquad(T\text{ en kelvins}) \]
Les trois lois du tableau n'en sont que des cas particuliers : il suffit de barrer la variable qui
ne change pas. C'est l'outil « couteau suisse » de ce thème.
\end{formule}

\section*{3. Deux exemples pour fixer les idées}

\textbf{(a) Compression isotherme (Boyle).} On enferme $\qty{24}{\litre}$ d'air à $\qty{1}{\atm}$ ; on
pousse le piston jusqu'à $\qty{3}{\atm}$ sans changer la température. Comme $T$ est constant :
\[ V_2=\frac{p_1V_1}{p_2}=\frac{1\times24}{3}=\qty{8}{\litre}. \]
On a triplé $p$, donc divisé $V$ par $3$ : le produit $pV=\qty{24}{\atm\litre}$ est conservé.

\textbf{(b) Chauffage isochore (bombe rigide).} Une bombe rigide contient un gaz à $\qty{2}{\atm}$ et
$\qty{27}{\degreeCelsius}$. On la chauffe à $\qty{127}{\degreeCelsius}$ ; $V$ ne change pas. D'abord les
kelvins : $T_1=\qty{300}{\kelvin}$, $T_2=\qty{400}{\kelvin}$. Puis $p_1/T_1=p_2/T_2$ :
\[ p_2=p_1\frac{T_2}{T_1}=2\times\frac{400}{300}=\qty{2.67}{\atm}. \]

\begin{methode}[La démarche, à appliquer à chaque exercice]
\textbf{1.} Lister l'état initial et l'état final ; \textbf{repérer la variable constante}. \quad
\textbf{2.} Convertir : $T$ en \textbf{kelvins}, et $p$, $V$ dans une unité \emph{cohérente des deux
côtés}. \quad
\textbf{3.} Choisir la loi (une variable figée $\to$ loi historique ; deux qui bougent $\to$ loi
combinée). \quad
\textbf{4.} Isoler l'inconnue \emph{avant} de mettre les nombres. \quad
\textbf{5.} Calculer, puis \textbf{vérifier le sens} : chauffer dilate, comprimer augmente $p$\dots
\end{methode}

\begin{idee}[Le réflexe « proportionnalité »]
Nul besoin de la valeur de $R$ ici : tout se joue en \textbf{rapports}. Si la température absolue est
multipliée par $1{,}2$ à pression constante, le volume est multiplié par $1{,}2$. Si la pression double
à température constante, le volume est divisé par $2$. Raisonner en facteurs multiplicatifs est souvent
plus rapide et plus sûr que de tout recalculer.
\end{idee}

\end{document}
